% =========================================================
%  Báo cáo Bài tập lớn — Matrix Chain Multiplication
%  Môn: CS112 — Phân tích và Thiết kế Thuật toán
%  Trường: Đại học Công nghệ Thông tin — ĐHQG TP.HCM (UIT)
%  GVHD: Huỳnh Thị Thanh Thương
%
%  Biên dịch: pdfLaTeX (mặc định) — không cần XeLaTeX/LuaLaTeX
%    pdflatex report.tex
%    pdflatex report.tex   (chạy 2 lần để cập nhật mục lục)
% =========================================================
\documentclass[12pt,a4paper]{article}

% --- ENCODING & FONT (pdfLaTeX — Unicode tiếng Việt qua inputenc + babel) ---
\usepackage[utf8]{inputenc}
\usepackage[T5]{fontenc}
\usepackage[vietnamese]{babel}
\usepackage{lmodern}    % font Latin Modern, hỗ trợ Unicode tốt
\usepackage{textcomp}

% --- LAYOUT ---
\usepackage[a4paper,margin=2.5cm,top=2.5cm,bottom=2.5cm]{geometry}
\usepackage{setspace}
\onehalfspacing
\usepackage{indentfirst}
\setlength{\parindent}{1.25em}
\setlength{\parskip}{0.4em}

% --- MATH ---
\usepackage{amsmath, amssymb, amsthm, mathtools}

% --- GRAPHICS / TABLES ---
\usepackage{graphicx}
\usepackage{array}
\usepackage{booktabs}
\usepackage{multirow}
\usepackage{tabularx}
\usepackage{caption}
\captionsetup{font=small,labelfont=bf}
\usepackage{float}

% --- COLORS & LINKS ---
\usepackage{xcolor}
\definecolor{primary}{HTML}{1d4ed8}
\definecolor{accent}{HTML}{92400e}
\definecolor{codebg}{HTML}{F8FAFC}
\definecolor{codeborder}{HTML}{E2E8F0}
\definecolor{codekw}{HTML}{7C3AED}
\definecolor{codestr}{HTML}{16A34A}
\definecolor{codecmt}{HTML}{64748B}
\usepackage[colorlinks=true,linkcolor=primary,urlcolor=primary,citecolor=primary]{hyperref}

% --- LISTINGS (CODE) ---
\usepackage{listings}
\lstdefinestyle{pythonstyle}{
  language=Python,
  basicstyle=\ttfamily\footnotesize,
  keywordstyle=\color{codekw}\bfseries,
  commentstyle=\color{codecmt}\itshape,
  stringstyle=\color{codestr},
  numbers=left, numberstyle=\tiny\color{codecmt},
  numbersep=8pt,
  backgroundcolor=\color{codebg},
  frame=single, framerule=0.4pt, rulecolor=\color{codeborder},
  breaklines=true, breakatwhitespace=true,
  showstringspaces=false,
  tabsize=4,
  captionpos=b,
  inputencoding=utf8,
  literate=
    {á}{{\'a}}1 {à}{{\`a}}1 {ả}{{\h{a}}}1 {ã}{{\~a}}1 {ạ}{{\d{a}}}1
    {â}{{\^a}}1 {ấ}{{\'\^a}}1 {ầ}{{\`\^a}}1 {ẩ}{{\h{\^a}}}1 {ẫ}{{\~\^a}}1 {ậ}{{\d{\^a}}}1
    {ă}{{\u a}}1 {ắ}{{\'\u a}}1 {ằ}{{\`\u a}}1 {ẳ}{{\h{\u a}}}1 {ẵ}{{\~\u a}}1 {ặ}{{\d{\u a}}}1
    {é}{{\'e}}1 {è}{{\`e}}1 {ẻ}{{\h e}}1 {ẽ}{{\~e}}1 {ẹ}{{\d e}}1
    {ê}{{\^e}}1 {ế}{{\'\^e}}1 {ề}{{\`\^e}}1 {ể}{{\h{\^e}}}1 {ễ}{{\~\^e}}1 {ệ}{{\d{\^e}}}1
    {í}{{\'i}}1 {ì}{{\`i}}1 {ỉ}{{\h i}}1 {ĩ}{{\~i}}1 {ị}{{\d i}}1
    {ó}{{\'o}}1 {ò}{{\`o}}1 {ỏ}{{\h o}}1 {õ}{{\~o}}1 {ọ}{{\d o}}1
    {ô}{{\^o}}1 {ố}{{\'\^o}}1 {ồ}{{\`\^o}}1 {ổ}{{\h{\^o}}}1 {ỗ}{{\~\^o}}1 {ộ}{{\d{\^o}}}1
    {ơ}{{\o }}1 {ớ}{{\'\o }}1 {ờ}{{\`\o }}1 {ở}{{\h{\o }}}1 {ỡ}{{\~\o }}1 {ợ}{{\d{\o }}}1
    {ú}{{\'u}}1 {ù}{{\`u}}1 {ủ}{{\h u}}1 {ũ}{{\~u}}1 {ụ}{{\d u}}1
    {ư}{{\u u}}1 {ứ}{{\'\u u}}1 {ừ}{{\`\u u}}1 {ử}{{\h{\u u}}}1 {ữ}{{\~\u u}}1 {ự}{{\d{\u u}}}1
    {ý}{{\'y}}1 {ỳ}{{\`y}}1 {ỷ}{{\h y}}1 {ỹ}{{\~y}}1 {ỵ}{{\d y}}1
    {đ}{{\dj}}1 {Đ}{{\DJ}}1
}
\lstset{style=pythonstyle}

% --- ALGORITHM ---
\usepackage{algorithm}
\usepackage{algpseudocode}
\floatname{algorithm}{Thuật toán}
\renewcommand{\algorithmicrequire}{\textbf{Đầu vào:}}
\renewcommand{\algorithmicensure}{\textbf{Đầu ra:}}

% --- HEADER & FOOTER ---
\usepackage{fancyhdr}
\pagestyle{fancy}
\fancyhf{}
\fancyhead[L]{\small\itshape Bài tập lớn — Matrix Chain Multiplication}
\fancyhead[R]{\small\itshape CS112 — UIT}
\fancyfoot[C]{\thepage}
\renewcommand{\headrulewidth}{0.4pt}

% --- MISC ---
\usepackage{enumitem}
\setlist{itemsep=2pt, topsep=4pt}

% =========================================================
%  TÀI LIỆU
% =========================================================
\begin{document}

% --- TRANG BÌA ---
\begin{titlepage}
  \centering
  \vspace*{1.5cm}

  {\large\textbf{ĐẠI HỌC QUỐC GIA THÀNH PHỐ HỒ CHÍ MINH}}\\[0.2cm]
  {\large\textbf{TRƯỜNG ĐẠI HỌC CÔNG NGHỆ THÔNG TIN}}\\[0.2cm]
  {\large\textbf{KHOA KHOA HỌC MÁY TÍNH}}

  \vspace{1.5cm}
  \includegraphics[width=0.22\textwidth]{example-image-a}\\[0.2cm]
  {\footnotesize \itshape (Thay placeholder bằng logo UIT nếu cần)}

  \vspace{1.5cm}
  {\Large\bfseries BÁO CÁO BÀI TẬP LỚN}\\[0.4cm]
  {\large Môn: CS112 — Phân tích và Thiết kế Thuật toán}

  \vspace{1.5cm}
  \rule{\linewidth}{0.5mm}\\[0.4cm]
  {\Huge\bfseries\color{primary} Matrix Chain Multiplication}\\[0.2cm]
  {\large\itshape Quy hoạch động Bottom-Up theo CLRS chương 15.2}\\[0.4cm]
  \rule{\linewidth}{0.5mm}

  \vspace{2cm}
  \begin{tabular}{ll}
    \textbf{Giảng viên hướng dẫn:} & ThS. Huỳnh Thị Thanh Thương \\[0.3em]
    \textbf{Sinh viên thực hiện:}  & \textit{[Họ và tên sinh viên]} \\[0.3em]
    \textbf{MSSV:}                 & \textit{[MSSV]} \\[0.3em]
    \textbf{Lớp:}                  & \textit{[Mã lớp]} \\[0.3em]
    \textbf{Học kỳ:}               & \textit{[Học kỳ — Năm học]} \\
  \end{tabular}

  \vfill
  {\large TP. Hồ Chí Minh, \today}
\end{titlepage}

% --- LỜI CAM ĐOAN ---
\section*{Lời cam đoan}
\addcontentsline{toc}{section}{Lời cam đoan}

Tôi cam đoan rằng báo cáo này là kết quả nghiên cứu và triển khai của cá nhân tôi. Mọi tham khảo từ các nguồn bên ngoài đều đã được trích dẫn rõ ràng. Phần triển khai bằng Python (gồm thuật toán lõi, kiểm thử, và giao diện web demo) được viết hoàn toàn trong khuôn khổ bài tập lớn của môn \textbf{CS112 — Phân tích và Thiết kế Thuật toán}.

\vspace{0.6em}
\hfill TP. Hồ Chí Minh, \today

\hfill \textit{Sinh viên thực hiện}

\vspace{2em}
\hfill \textit{[Ký và ghi rõ họ tên]}

\newpage

% --- TÓM TẮT ---
\section*{Tóm tắt}
\addcontentsline{toc}{section}{Tóm tắt}

Báo cáo trình bày quá trình phân tích, thiết kế và triển khai bài toán \textbf{Matrix Chain Multiplication (MCM)} bằng phương pháp \textbf{Quy hoạch động Bottom-Up} theo mã giả CLRS chương~15.2. Hệ thống được xây dựng bằng \textbf{Python 3.10+} (chỉ sử dụng thư viện chuẩn cho phần lõi), bám đúng từng dòng mã giả của giáo trình bằng kỹ thuật \emph{1-based padding} để code đọc gần với toán học của slide bài giảng.

Phần triển khai gồm bốn module độc lập (\texttt{validator}, \texttt{solver}, \texttt{formatter}, \texttt{cli}) với 43 test case (gồm 11 \emph{property-based test} kèm cross-check brute-force) đảm bảo tính đúng đắn về mặt toán học và định dạng. Ngoài giao diện dòng lệnh (CLI), một \textbf{ứng dụng web demo} viết bằng \textbf{Flask} được kèm theo để minh hoạ trực quan: bảng $m$, bảng $s$, các bước tính toán theo từng độ dài chuỗi con, cây phân tách (\emph{parenthesization tree}) và biểu thức ngoặc tối ưu.

Toàn bộ kết quả đối chiếu khớp với các ví dụ kinh điển trong giáo trình CLRS và slide bài giảng môn học.

\textbf{Từ khoá:} Matrix Chain Multiplication, Quy hoạch động, Bottom-Up, CLRS, Property-Based Testing, Flask.

\newpage

% --- MỤC LỤC ---
\tableofcontents
\newpage

% =========================================================
\section{Giới thiệu}
% =========================================================

\subsection{Bối cảnh}
Phép nhân hai ma trận $A$ kích thước $p \times q$ với $B$ kích thước $q \times r$ cần đúng $p \cdot q \cdot r$ phép nhân vô hướng. Khi nhân chuỗi $n$ ma trận $A_1 A_2 \cdots A_n$, do phép nhân ma trận có tính \emph{kết hợp}, có nhiều cách đặt ngoặc khác nhau. Mỗi cách cho kết quả tích giống nhau nhưng \emph{tổng số phép nhân vô hướng} có thể khác nhau \emph{rất nhiều}.

Ví dụ với $A_1 (10 \times 100), A_2 (100 \times 5), A_3 (5 \times 50)$:
\begin{itemize}
  \item $((A_1 A_2) A_3)$: $10\cdot100\cdot5 + 10\cdot5\cdot50 = 5000 + 2500 = \mathbf{7500}$ phép nhân.
  \item $(A_1 (A_2 A_3))$: $100\cdot5\cdot50 + 10\cdot100\cdot50 = 25000 + 50000 = \mathbf{75000}$ phép nhân.
\end{itemize}

Cách đặt ngoặc đầu tiên \emph{nhanh hơn 10 lần}. Bài toán MCM đặt ra: cho trước kích thước của $n$ ma trận, tìm cách đặt ngoặc tối ưu để \textbf{tối thiểu hoá} tổng số phép nhân vô hướng.

\subsection{Mục tiêu của bài tập lớn}
\begin{enumerate}
  \item Phát biểu bài toán MCM một cách hình thức.
  \item Phân tích cấu trúc con tối ưu, thiết kế công thức truy hồi và thuật toán Bottom-Up bằng quy hoạch động.
  \item Triển khai thuật toán bằng Python theo đúng mã giả CLRS, có kiểm thử tự động (unit + property-based).
  \item Xây dựng giao diện web demo (Flask) để minh hoạ trực quan các bước, bảng $m$/$s$ và cây phân tách.
  \item Kiểm chứng độ phức tạp $\Theta(n^3)$ qua test hiệu năng và đối chiếu với các ví dụ trong giáo trình.
\end{enumerate}

\subsection{Đóng góp chính}
\begin{itemize}
  \item Cài đặt bám sát từng dòng mã giả CLRS bằng \emph{1-based padding}, giúp ánh xạ trực tiếp giữa code Python và mã giả không cần dịch chỉ số.
  \item Bộ kiểm thử 43 test gồm 11 \emph{property-based test} thuần thư viện chuẩn (không phụ thuộc \texttt{hypothesis}), kèm cross-check brute-force liệt kê \emph{toàn bộ} parenthesization với $n \le 6$.
  \item Giao diện web demo trực quan với ba tab: bảng $m$/$s$, các bước tính toán có nút \emph{Trước/Sau/Tự động}, cây phân tách của biểu thức tối ưu.
\end{itemize}

% =========================================================
\section{Phát biểu bài toán}
% =========================================================

\subsection{Đầu vào và đầu ra}
\textbf{Đầu vào.} Một mảng số nguyên dương $p = [p_0, p_1, \ldots, p_n]$ gồm $n+1$ phần tử, mô tả chuỗi $n$ ma trận:
\[
A_i \in \mathbb{R}^{p_{i-1} \times p_i}, \qquad 1 \le i \le n.
\]

\textbf{Đầu ra.}
\begin{itemize}
  \item $m[1, n]$ — số phép nhân vô hướng tối thiểu.
  \item Bảng $s$ ghi vị trí tách tối ưu, từ đó truy vết được \emph{Optimal Parenthesization}.
\end{itemize}

\subsection{Định nghĩa hình thức}
Đặt $m[i, j]$ là số phép nhân vô hướng tối thiểu để tính tích $A_i A_{i+1} \cdots A_j$ với $1 \le i \le j \le n$. Khi đó:
\begin{equation}\label{eq:mcm}
m[i, j] =
\begin{cases}
0, & \text{nếu } i = j, \\[6pt]
\displaystyle \min_{i \le k < j} \Big\{ m[i, k] + m[k+1, j] + p_{i-1}\, p_k\, p_j \Big\}, & \text{nếu } i < j.
\end{cases}
\end{equation}

Đáp số cuối là $m[1, n]$.

\subsection{Cấu trúc con tối ưu}
Một parenthesization tối ưu của $A_i \cdots A_j$ phải tách tại một vị trí $k$ nào đó ($i \le k < j$) thành hai chuỗi con $A_i \cdots A_k$ và $A_{k+1} \cdots A_j$. Hai chuỗi con đó cũng phải được parenthesize tối ưu — đây là tính chất \emph{tối ưu lồng tối ưu} điển hình của quy hoạch động.

\subsection{Phân tích trùng lặp con}
Số chuỗi con phân biệt $(i, j)$ là $\binom{n}{2} + n = \Theta(n^2)$. Nếu giải đệ quy không nhớ, mỗi chuỗi con sẽ được tính lại nhiều lần dẫn đến độ phức tạp \emph{số Catalan} $\Omega(4^n / n^{3/2})$. Do đó cần \textbf{quy hoạch động} để mỗi chuỗi con chỉ tính một lần.

% =========================================================
\section{Thuật toán}
% =========================================================

\subsection{Mã giả \texttt{MATRIX-CHAIN-ORDER}}

\begin{algorithm}[H]
\caption{\textsc{Matrix-Chain-Order}$(p)$}
\begin{algorithmic}[1]
\Require Mảng $p[0..n]$ gồm $n+1$ số nguyên dương.
\Ensure Bảng $m[1..n, 1..n]$ và bảng $s[1..n-1, 2..n]$.
\State $n \gets p.\text{length} - 1$
\State Cấp phát hai bảng $m$ và $s$
\For{$i \gets 1$ \textbf{to} $n$}
  \State $m[i, i] \gets 0$
\EndFor
\For{$\ell \gets 2$ \textbf{to} $n$} \Comment{$\ell$ là độ dài chuỗi con}
  \For{$i \gets 1$ \textbf{to} $n - \ell + 1$}
    \State $j \gets i + \ell - 1$
    \State $m[i, j] \gets +\infty$
    \For{$k \gets i$ \textbf{to} $j - 1$}
      \State $q \gets m[i, k] + m[k+1, j] + p_{i-1} \cdot p_k \cdot p_j$
      \If{$q < m[i, j]$}
        \State $m[i, j] \gets q$
        \State $s[i, j] \gets k$
      \EndIf
    \EndFor
  \EndFor
\EndFor
\State \Return $(m, s)$
\end{algorithmic}
\end{algorithm}

\subsection{Mã giả \texttt{PRINT-OPTIMAL-PARENS}}

\begin{algorithm}[H]
\caption{\textsc{Print-Optimal-Parens}$(s, i, j)$}
\begin{algorithmic}[1]
\If{$i = j$}
  \State In ``$A_i$''
\Else
  \State In ``$($''
  \State \textsc{Print-Optimal-Parens}$(s, i, s[i, j])$
  \State \textsc{Print-Optimal-Parens}$(s, s[i, j] + 1, j)$
  \State In ``$)$''
\EndIf
\end{algorithmic}
\end{algorithm}

\subsection{Phân tích độ phức tạp}

\begin{itemize}
  \item \textbf{Thời gian.} Ba vòng lặp lồng nhau theo $\ell, i, k$, mỗi cặp $(i, j, k)$ thực hiện $O(1)$ phép tính. Tổng số phép gán $q$ là
  \[
  \sum_{\ell=2}^{n} \sum_{i=1}^{n-\ell+1} (\ell - 1) \;=\; \binom{n}{3} + \binom{n}{2} \;=\; \Theta(n^3).
  \]
  \item \textbf{Không gian.} Hai bảng $m$ và $s$ kích thước $\Theta(n^2)$, đệ quy in parens có chiều sâu tối đa $O(n)$. Tổng cộng $\Theta(n^2)$.
  \item \textbf{So sánh với đệ quy không nhớ.} Đệ quy thuần có độ phức tạp $\Omega(4^n / n^{3/2})$ (số Catalan), \emph{vô vọng} với $n$ lớn. Quy hoạch động cải thiện về $\Theta(n^3)$.
\end{itemize}

\subsection{Ví dụ chạy tay (CLRS kinh điển)}

Cho $p = [30, 35, 15, 5, 10, 20, 25]$, ta có 6 ma trận:
\[
A_1: 30{\times}35,\; A_2: 35{\times}15,\; A_3: 15{\times}5,\; A_4: 5{\times}10,\; A_5: 10{\times}20,\; A_6: 20{\times}25.
\]

Sau khi chạy thuật toán, bảng $m$ và $s$ thu được như Bảng~\ref{tab:m-clrs} và \ref{tab:s-clrs}. Đáp số cuối:
\[
m[1, 6] = 15{,}125, \qquad \text{Optimal Parens} = ((A_1(A_2A_3))((A_4A_5)A_6)).
\]

\begin{table}[H]
\centering
\caption{Bảng $m$ cho $p = [30, 35, 15, 5, 10, 20, 25]$.}
\label{tab:m-clrs}
\small
\begin{tabular}{c|rrrrrr}
\toprule
$i \backslash j$ & $j=1$ & $j=2$ & $j=3$ & $j=4$ & $j=5$ & $j=6$ \\
\midrule
$i=1$ & $0$ & $15750$ & $7875$ & $9375$ & $11875$ & $\mathbf{15125}$ \\
$i=2$ & $-$ & $0$     & $2625$ & $4375$ & $7125$  & $10500$ \\
$i=3$ & $-$ & $-$     & $0$    & $750$  & $2500$  & $5375$  \\
$i=4$ & $-$ & $-$     & $-$    & $0$    & $1000$  & $3500$  \\
$i=5$ & $-$ & $-$     & $-$    & $-$    & $0$     & $5000$  \\
$i=6$ & $-$ & $-$     & $-$    & $-$    & $-$     & $0$     \\
\bottomrule
\end{tabular}
\end{table}

\begin{table}[H]
\centering
\caption{Bảng $s$ — vị trí tách tối ưu.}
\label{tab:s-clrs}
\small
\begin{tabular}{c|rrrrr}
\toprule
$i \backslash j$ & $j=2$ & $j=3$ & $j=4$ & $j=5$ & $j=6$ \\
\midrule
$i=1$ & $1$ & $1$ & $3$ & $3$ & $\mathbf{3}$ \\
$i=2$ & $-$ & $2$ & $3$ & $3$ & $3$ \\
$i=3$ & $-$ & $-$ & $3$ & $3$ & $3$ \\
$i=4$ & $-$ & $-$ & $-$ & $4$ & $5$ \\
$i=5$ & $-$ & $-$ & $-$ & $-$ & $5$ \\
\bottomrule
\end{tabular}
\end{table}

% =========================================================
\section{Triển khai}
% =========================================================

\subsection{Lựa chọn công nghệ}
\begin{itemize}
  \item \textbf{Ngôn ngữ:} Python 3.10+ (chỉ thư viện chuẩn cho phần lõi). Cú pháp Python đọc gần với mã giả, thuận tiện cho mục đích giảng dạy.
  \item \textbf{Quy ước chỉ số:} \emph{1-based padding} — bảng $m, s$ kích thước $(n+1) \times (n+1)$, ô index $0$ không dùng. Nhờ vậy code Python truy cập \texttt{m[i][j]}, \texttt{s[i][j]} \emph{mirror đúng từng dòng} mã giả CLRS, tránh phải dịch \texttt{[i-1]} ở mọi nơi.
  \item \textbf{Sentinel cho $\infty$:} \texttt{math.inf} (Python) trong vòng so sánh, sau khi vòng $k$ kết thúc giá trị chắc chắn là số nguyên do tồn tại ít nhất một $q < +\infty$.
  \item \textbf{Test framework:} \texttt{pytest} (dev-dependency duy nhất). Property-based test viết thuần bằng \texttt{random.Random(seed)} cố định để đảm bảo tái lập mà không phụ thuộc \texttt{hypothesis}.
\end{itemize}

\subsection{Cấu trúc dự án}
\begin{lstlisting}[language=,basicstyle=\ttfamily\small,frame=single,backgroundcolor=\color{codebg}]
matrix-chain/
├── src/mcm/
│   ├── __init__.py       # re-export public API
│   ├── solver.py         # MATRIX-CHAIN-ORDER, PRINT-OPTIMAL-PARENS, MCMSolver
│   ├── validator.py      # validate_dimension_array
│   ├── formatter.py      # DisplayMode, format_table, format_trace
│   ├── cli.py            # entry point CLI (argparse)
│   └── __main__.py       # python -m mcm
├── tests/
│   ├── test_validator.py # 12 tests
│   ├── test_solver.py    # 16 tests + brute-force cross-check + perf
│   ├── test_formatter.py # 12 tests
│   └── test_cli.py       # 3 smoke tests
├── templates/index.html  # giao diện web (UI demo)
├── static/style.css      # CSS cho UI
├── web_app.py            # Flask app
├── pyproject.toml
└── README.md
\end{lstlisting}

\subsection{Module \texttt{validator}}
Hàm \texttt{validate\_dimension\_array(p) -> int} kiểm tra ba điều kiện theo thứ tự: độ dài $\ge 2$, kiểu \texttt{int} (loại trừ \texttt{bool} — vốn là subclass của \texttt{int} trong Python), và giá trị dương. Khi hợp lệ, trả về $n = \texttt{len}(p) - 1$.

Thông điệp lỗi tiếng Việt được giữ \emph{nguyên văn} theo yêu cầu trong tài liệu môn học, ví dụ: ``\texttt{Mọi phần tử của Dimension\_Array phải là số nguyên dương}''.

\subsection{Module \texttt{solver} — phần lõi}

\begin{lstlisting}[caption={\texttt{matrix\_chain\_order} — bám đúng từng dòng mã giả CLRS.},label={lst:mco}]
import math

def matrix_chain_order(p):
    n = len(p) - 1                                    # dòng 1
    m = [[0] * (n + 1) for _ in range(n + 1)]         # dòng 2 (1-based padding)
    s = [[0] * (n + 1) for _ in range(n + 1)]
    for i in range(1, n + 1):                         # dòng 3-4
        m[i][i] = 0

    for length in range(2, n + 1):                    # dòng 5: l = 2..n
        for i in range(1, n - length + 2):            # dòng 6
            j = i + length - 1                        # dòng 7
            m[i][j] = math.inf                        # dòng 8
            for k in range(i, j):                     # dòng 9
                q = m[i][k] + m[k+1][j] + p[i-1]*p[k]*p[j]   # dòng 10
                if q < m[i][j]:                       # dòng 11
                    m[i][j] = q                       # dòng 12
                    s[i][j] = k                       # dòng 13
    return m, s
\end{lstlisting}

Hàm \texttt{print\_optimal\_parens(s, i, j)} dùng \texttt{StringIO} làm accumulator để giữ cấu trúc đệ quy ``in \texttt{`(`}, gọi đệ quy, in \texttt{`)`}'' giống mã giả, đồng thời trả về kiểu \texttt{str} để dễ unit-test.

Lớp \texttt{MCMSolver} đóng vai trò \emph{facade}, gói validator $\to$ solver $\to$ print parens $\to$ tự kiểm tra nhất quán nội bộ trong một call duy nhất. Sau khi giải xong, mọi truy vấn (\texttt{cost}, \texttt{parens}, \texttt{m}, \texttt{s}, \texttt{trace}) đều có thể lấy nhanh.

\subsection{Module \texttt{formatter}}
\begin{itemize}
  \item Enum \texttt{DisplayMode} với ba giá trị: \texttt{NONE} / \texttt{PARTIAL\_TABLES} / \texttt{FULL\_TABLES}.
  \item \texttt{format\_table} căn lề phải, ô không dùng (ô $i > j$ với bảng $m$, ô $i \ge j$ với bảng $s$) hiển thị dấu \texttt{`-`}. Độ rộng cột thống nhất bằng $\max\{\texttt{len(str(value))}\} + 2$.
  \item \texttt{format\_trace} ghép bảng + dòng kết quả thành một \emph{Trace\_Output} với bốn nhãn tiếng Việt nguyên văn.
\end{itemize}

\subsection{Module \texttt{cli}}
Sử dụng \texttt{argparse}: positional \texttt{dims} (đọc từ \texttt{stdin} nếu rỗng), tham số \texttt{--mode} (\{\texttt{none, partial, full}\}, mặc định \texttt{full}) và \texttt{--table} (\{\texttt{m, s}\}, áp dụng khi \texttt{--mode=partial}).

Quy ước mã thoát: $0$ thành công, $1$ \texttt{ValueError} (input sai), $2$ \texttt{RuntimeError} hoặc lỗi khác (ví dụ vi phạm tính nhất quán nội bộ). Khi $n > 500$, CLI gọi \texttt{sys.setrecursionlimit} trước khi truy vết bảng $s$ để tránh tràn stack.

% =========================================================
\section{Kiểm thử và đảm bảo chất lượng}
% =========================================================

\subsection{Tổng quan}
Bộ test có \textbf{43 test case} chia thành 4 nhóm: \texttt{validator}, \texttt{solver}, \texttt{formatter}, \texttt{cli}. Toàn bộ 43 test \emph{pass} trong $\approx 0.30$ giây.

\subsection{Property-Based Testing (PBT)}
Property-Based Testing là kỹ thuật kiểm thử bằng cách phát biểu các \emph{tính chất bất biến} mà chương trình phải thoả mãn với \emph{mọi} đầu vào hợp lệ, sau đó kiểm tra bằng cách sinh ngẫu nhiên một số lượng lớn input. Do yêu cầu chỉ dùng thư viện chuẩn, ở đây tôi không dùng \texttt{hypothesis} mà dùng \texttt{random.Random(SEED)} cố định, mỗi property một offset trên seed cơ sở $20251128$ để \emph{tái lập} 100\%.

Bảng~\ref{tab:pbt} liệt kê 11 property đã được kiểm.

\begin{table}[H]
\centering
\caption{11 property-based test với cross-check brute-force.}
\label{tab:pbt}
\renewcommand{\arraystretch}{1.18}
\small
\begin{tabularx}{\linewidth}{c >{\raggedright\arraybackslash}X c}
\toprule
\textbf{\#} & \textbf{Phát biểu tính chất} & \textbf{$n_{\max}$} \\
\midrule
1  & Validator trả về $\texttt{len}(p) - 1$                                          & 8 \\
2  & Validator từ chối phần tử $\le 0$                                                & 8 \\
3  & $m[i][i] = 0$ với mọi $i$ (đường chéo)                                            & 8 \\
4  & Tự nhất quán: $s[i][j] \in [i, j-1]$, recurrence cực tiểu, $m[i][j]$ là min       & 8 \\
5  & \textbf{Cross-check brute-force:} $m[1][n]$ khớp min trên \emph{mọi} parenthesize  & 6 \\
6  & API \texttt{cost} khớp $m[1][n]$                                                  & 8 \\
7  & Cấu trúc parens hợp lệ (ngoặc cân bằng, token $A_k$ thứ tự đúng)                  & 8 \\
8  & Ô không sử dụng trong bảng hiển thị \texttt{`-`}                                  & 8 \\
9  & Cột bảng căn lề nhất quán                                                         & 8 \\
10 & Trace\_Output \texttt{FULL\_TABLES} chứa đủ 4 nhãn tiếng Việt                     & 8 \\
11 & Cost in dạng số nguyên thuần (regex \texttt{\textbackslash d+})                   & 8 \\
\bottomrule
\end{tabularx}
\end{table}

Mỗi property test chạy 100 iteration, tổng cộng $\approx 1100$ test case sinh tự động. Property~5 đặc biệt quan trọng: nó so sánh $m[1][n]$ của thuật toán quy hoạch động với chi phí tối thiểu khi liệt kê \emph{toàn bộ} parenthesization (số Catalan $C_{n-1}$), do đó cross-check tính \emph{tối ưu Bellman} của bảng $m$.

\subsection{Unit test 4 ví dụ bắt buộc}

\begin{table}[H]
\centering
\caption{Bốn ví dụ bắt buộc theo tài liệu môn học.}
\small
\begin{tabular}{l c c}
\toprule
\textbf{Dimension Array $p$} & $\boldsymbol{m[1, n]}$ & \textbf{Optimal Parens} \\
\midrule
$[3, 10, 7, 2]$               & $200$    & $(A_1(A_2A_3))$ \\
$[10, 20, 50, 1, 100]$        & $2200$   & $((A_1(A_2A_3))A_4)$ \\
$[30, 35, 15, 5, 10, 20, 25]$ & $15125$  & $((A_1(A_2A_3))((A_4A_5)A_6))$ \\
$[5, 7]$                      & $0$      & $A_1$ \\
\bottomrule
\end{tabular}
\end{table}

\subsection{Test hiệu năng}
\begin{itemize}
  \item Với $n = 100$: hoàn thành trong $\approx 0.05$ giây trên máy tham chiếu (yêu cầu giáo trình $< 1\,\text{s}$).
  \item Với $n = 150$: $\approx 0.18$ giây — không reject đầu vào hợp lệ (graceful degradation).
\end{itemize}

% =========================================================
\section{Giao diện web demo}
% =========================================================

\subsection{Kiến trúc tổng thể}
Ứng dụng web demo viết bằng \textbf{Flask}, sử dụng kiến trúc đơn giản:
\begin{itemize}
  \item \texttt{web\_app.py}: Flask server với 2 endpoint.
    \begin{itemize}[noitemsep]
      \item \texttt{GET /}: render template \texttt{index.html}.
      \item \texttt{POST /api/solve}: nhận JSON \texttt{\{dims: [...]\}}, gọi \texttt{MCMSolver}, trả về JSON gồm \texttt{cost}, \texttt{parens}, bảng $m$, bảng $s$, và danh sách các bước tính toán.
    \end{itemize}
  \item \texttt{templates/index.html}: giao diện chính (HTML + Vanilla JS, không dùng framework JS).
  \item \texttt{static/style.css}: thiết kế responsive, palette màu xanh navy + accent vàng cho ô tối ưu.
\end{itemize}

\subsection{Tính năng giao diện}
\begin{enumerate}
  \item \textbf{Nhập Dimension Array} dưới dạng văn bản (hỗ trợ phân cách bằng dấu phẩy hoặc khoảng trắng), kèm 4 ví dụ mẫu bấm nhanh.
  \item \textbf{Tab \emph{Bảng $m$ \& $s$}:} hiển thị hai bảng đầy đủ, ô tối ưu $m[1, n]$ tô vàng, ô không dùng để dấu \texttt{`-`}.
  \item \textbf{Tab \emph{Các bước}:} accordion từng bước theo độ dài chuỗi con $\ell$, cho từng cặp $(i, j)$ liệt kê \emph{mọi} candidate $k$ với công thức đầy đủ và đánh dấu candidate tối ưu. Có nút \emph{Trước / Sau / Tự động} để xem từng bước, đồng thời ô đang tính được highlight trên bảng.
  \item \textbf{Tab \emph{Cây đóng ngoặc}:} render cây phân tách theo bảng $s$, giúp sinh viên hình dung trực quan thứ tự nhân ma trận. Bên dưới là phần \emph{truy vết} dạng văn bản (textual trace).
\end{enumerate}

\subsection{Mẫu kết quả}
Với input $p = [30, 35, 15, 5, 10, 20, 25]$, giao diện hiển thị:
\begin{itemize}
  \item Chi phí tối thiểu: \textbf{15{,}125} (tô vàng đậm).
  \item Đóng ngoặc tối ưu: \texttt{((A1(A2A3))((A4A5)A6))}.
  \item 16 bước tính toán (1 bước khởi tạo + 15 cặp $(i, j)$ với $i < j$).
  \item Cây phân tách 3 tầng theo $s[1, 6] = 3$, $s[1, 3] = 1$, $s[4, 6] = 5$.
\end{itemize}

\subsection{Cách chạy}
\begin{lstlisting}[language=,basicstyle=\ttfamily\small,frame=single,backgroundcolor=\color{codebg}]
# Lần đầu: tạo và kích hoạt môi trường ảo
python -m venv venv
.\venv\Scripts\activate   (Windows) hoặc source venv/bin/activate (Linux/macOS)

# Cài đặt
pip install -e ".[dev]" flask

# Chạy
python web_app.py
# Mở http://127.0.0.1:5000 trên trình duyệt
\end{lstlisting}

% =========================================================
\section{Kết quả thực nghiệm}
% =========================================================

\subsection{Đầu ra mẫu của CLI}
Lệnh \texttt{python -m mcm 30 35 15 5 10 20 25 --mode=full} cho output:

\begin{lstlisting}[language=,basicstyle=\ttfamily\footnotesize,frame=single,backgroundcolor=\color{codebg}]
Bảng m:
            j=1     j=2     j=3     j=4     j=5     j=6
    i=1       0   15750    7875    9375   11875   15125
    i=2       -       0    2625    4375    7125   10500
    i=3       -       -       0     750    2500    5375
    i=4       -       -       -       0    1000    3500
    i=5       -       -       -       -       0    5000
    i=6       -       -       -       -       -       0
Bảng s:
        j=2   j=3   j=4   j=5   j=6
  i=1     1     1     3     3     3
  i=2     -     2     3     3     3
  i=3     -     -     3     3     3
  i=4     -     -     -     4     5
  i=5     -     -     -     -     5
Cách đóng mở ngoặc tối ưu: ((A1(A2A3))((A4A5)A6))
Số phép nhân vô hướng tối thiểu: 15125
\end{lstlisting}

\subsection{Đối chiếu với CLRS}
Cả hai bảng và biểu thức ngoặc đều khớp \emph{tuyệt đối} với Hình 15.5 của CLRS chương 15.2 \cite{clrs}.

% =========================================================
\section{Kết luận}
% =========================================================

\subsection{Kết quả đạt được}
Báo cáo đã trình bày trọn vẹn quá trình giải bài toán Matrix Chain Multiplication: từ phát biểu hình thức, công thức truy hồi, mã giả, đến triển khai bằng Python bám sát từng dòng mã giả CLRS. Hệ thống có:
\begin{itemize}
  \item Bốn module Python tổ chức rõ ràng (validator, solver, formatter, cli).
  \item Bộ kiểm thử 43 test case (gồm 11 property-based test với cross-check brute-force liệt kê toàn bộ Catalan parenthesization).
  \item Giao diện CLI với ba chế độ hiển thị bảng (\texttt{none}, \texttt{partial}, \texttt{full}).
  \item Giao diện web demo (Flask) trực quan với ba tab: bảng, các bước, cây phân tách.
\end{itemize}

\subsection{Hạn chế}
\begin{itemize}
  \item Mới triển khai phiên bản Bottom-Up; chưa làm phiên bản top-down memoized để so sánh.
  \item Cây phân tách chỉ render dạng cây HTML thuần, chưa có animation từng bước trên bảng.
  \item Web app dùng Flask development server, chưa đóng gói WSGI server cho deployment thật.
\end{itemize}

\subsection{Hướng mở rộng}
\begin{itemize}
  \item Thêm phiên bản top-down (memoization) và đo benchmark so với Bottom-Up.
  \item Tích hợp thuật toán Hu--Shing $O(n \log n)$ làm tham chiếu nâng cao cho sinh viên hứng thú.
  \item Triển khai cây phân tách dạng SVG có animation, đồng bộ với từng bước trên bảng $m$.
  \item Triển khai phép nhân ma trận thực tế (không chỉ tính chi phí) để minh hoạ end-to-end.
\end{itemize}

% =========================================================
\begin{thebibliography}{9}
\addcontentsline{toc}{section}{Tài liệu tham khảo}

\bibitem{clrs}
T. H. Cormen, C. E. Leiserson, R. L. Rivest, and C. Stein.
\newblock \emph{Introduction to Algorithms}, 4th edition.
\newblock MIT Press, 2022.

\bibitem{slide}
Huỳnh Thị Thanh Thương.
\newblock \emph{Slide bài giảng môn ``Phân tích và Thiết kế Thuật toán''}.
\newblock Trường Đại học Công nghệ Thông tin (UIT), 2024--2025.

\bibitem{python}
Python Software Foundation.
\newblock \emph{Python 3.10 Documentation}.
\newblock \url{https://docs.python.org/3.10/}.

\bibitem{flask}
Pallets Projects.
\newblock \emph{Flask Documentation}.
\newblock \url{https://flask.palletsprojects.com/}.

\bibitem{pbt}
J. Hughes.
\newblock QuickCheck: Property-based testing for Haskell.
\newblock In \emph{ICFP}, 2000.

\end{thebibliography}

\end{document}
